In cyber-physical systems (CPS), explainable AI (XAI) helps operators act on intrusion alarms only if it identifies the compromised part of the process before it becomes unsafe. We introduce deadline-conditioned metrics that combine detection delay, explainer latency and the correctness of the delivered explanation, and evaluate explanations at the level of control loops. We curate loop-level ground truth for all 50 attacks of the HAI~21.03 testbed, cross-checked against an independent transcription, and release PCL-Sim, a closed-loop simulator whose \SimNattacks{} attacks carry physical hazard deadlines. We compare 13 attribution configurations, including exact loop Shapley values (\LX), a physics-informed loop attribution (PI-CAX) and temporal evidence accumulation (\TEA), across four detector families with leakage-free protocols. Score attributions of every kind degrade after the alarm as attacks propagate through coupled loops (HAI: \LX{} falls from \lochaiLXalarmZerosinglehrOne{} to \lochaiLXalarmSixZerosinglehrOne{} top-1 accuracy within 60\,s); \TEA{} recovers part of this loss. PI-CAX, which separates the persistent departure of the manipulated loop from the oscillatory compensation of the others, is significantly more accurate than every other configuration: on HAI, \lochaiPICAXalarmSixZeroteahrOne{} at 60\,s, and \dclhaiPICAXdclfixedSixZerotea{} versus at most \dclhaiIGdclfixedSixZerotea{} for correct localization within 60\,s of onset. Its advantage comes almost entirely from this component; it needs recalibration when sensor noise changes. On hazardous simulated attacks, \intimeFailUndeliveredPct\,\% of deadline misses are due to late or missed detection.
